% From Meal Prep.tex. Appliance: Instant Pot (5.7 L)
\begin{recipebox}{Budget Lentil Coconut Curry}{6}{Prep 10 min, cook 12 min pressure + 10 min natural release}
\subsection*{Why This Works}
\begin{itemize}
    \item Very low cost per serving
    \item High protein and fibre from lentils
    \item Filling but easy to digest
    \item Excellent meat-free contrast to heavier stews
    \item Ideal for lunches and reheats well
    \item Frozen vegetables keep cost low and eliminate waste
\end{itemize}

\subsection*{Ingredients}
\begin{itemize}
    \item Dried red lentils — 2$\tfrac{1}{2}$ cups, rinsed
    \item Onion — 1 large, diced
    \item Garlic — 3 cloves, minced
    \item Curry powder — 2 tbsp
    \item Ground cumin — 1 tbsp
    \item Canned tomatoes — 2 $\times$ 400 g
    \item Light coconut milk — 1 $\times$ 400 ml
    \item Water — 2 cups (500 ml)
    \item Salt — 1$\tfrac{1}{2}$ tsp
    \item Frozen vegetables — 1--2 cups (see options below)
\end{itemize}

\subsubsection*{Recommended Frozen Vegetable Options}
\begin{itemize}
    \item Pam's Mixed Vegetables (peas, corn, carrot, beans) — best value and texture
    \item Broccoli \& cauliflower mix — higher fibre, slightly firmer texture
    \item Garden peas — optional for added bulk and mild sweetness
    \item Frozen spinach — optional; use as a nutrient boost rather than bulk
\end{itemize}

\subsection*{Method}
\begin{enumerate}
    \item Add the lentils, onion, garlic, curry powder, and ground cumin to the Instant Pot.
    \item Add the canned tomatoes, coconut milk, water, and salt.
    \item Stir well to combine, ensuring lentils are fully submerged.
    \item Secure the lid with the valve sealed and cook on \textbf{High Pressure} for 12 minutes.
    \item Allow a \textbf{Natural Pressure Release} for 10 minutes, then vent remaining pressure.
    \item Open the lid and stir thoroughly; the lentils will thicken rapidly.
    \item Stir in frozen vegetables.
    \item Select \textbf{Sauté} mode and cook for 5--7 minutes until vegetables are heated through.
    \item Adjust thickness with additional water if required.
\end{enumerate}

\subsection*{Portioning \& Storage}
\begin{itemize}
    \item All 6 servings: refrigerate
    \item Best consumed within 4 days
    \item Thickens further after chilling; thin with water when reheating
\end{itemize}

\subsection*{What to Prep Alongside}

\paragraph{Rice}
\begin{itemize}
    \item Dry rice — 2$\tfrac{1}{2}$ cups
    \item Yield — approximately 7 cups cooked
\end{itemize}

\paragraph{Vegetables}
\begin{itemize}
    \item No additional vegetables required; curry is designed as a one-pot meal
    \item Optional: extra steamed frozen veg for larger portions
\end{itemize}

\subsection*{Serving Ideas}
\begin{itemize}
    \item \textbf{Lunch:} Lentil coconut curry, $\tfrac{3}{4}$ cup rice
    \item \textbf{Dinner:} Lentil coconut curry, $\tfrac{1}{2}$ cup rice, extra vegetables if desired
\end{itemize}

\subsection*{Notes}
\begin{itemize}
    \item Frozen pumpkin or diced potatoes (2--3 cups) may be added for extra bulk without changing cook time.
    \item Full-fat coconut milk can be used for richer flavour if budget allows.
    \item This curry freezes well; texture will be thicker after thawing.
\end{itemize}
\end{recipebox}
